\chapter{Introduction: The Crisis of Stochasticity}
\section{The Fragility of AI-Assisted Process Mining}
Modern process discovery has increasingly turned toward machine learning (ML) and reinforcement learning (RL) to handle the complexity of real-world event logs. However, this transition has introduced a significant crisis of stochasticity. Standard RL implementations rely heavily on dynamically allocated memory, non-deterministic bucket-probing hash maps, and floating-point approximations. These architectural choices lead to thread-level jitter and unstable execution trajectories. 

When an algorithm's output relies on the transient state of the host machine's heap or the CPU's branch predictor, the resulting process models become non-reproducible. In high-stakes environments—such as medical workflows, financial compliance, or legal dispatch systems—the ability to precisely reproduce a process model across different hardware platforms is not merely a convenience; it is a fundamental requirement for technical auditability and legal defensibility.

\section{The ``Cracked Contest'' and the Overfitting Accusation}
The process mining community frequently evaluates algorithms based on precision, recall, and classification accuracy across standardized datasets (e.g., the PDC-2025 suite). Achieving 100\% accuracy is often met with immediate skepticism, primarily due to the ubiquitous problem of \textit{overfitting}. Traditional algorithms that achieve perfect scores usually do so by memorizing trace sequences, resulting in overly complex "flower nets" or highly fragmented models that fail to generalize.

To claim 100\% accuracy legitimately, a system must prove that its models are not just accurate, but structurally minimal and logically sound. It must prove that it learned the underlying control-flow constraints rather than merely memorizing the training log.

\section{The DTEAM Vision: From Heuristics to Determinism}
The Deterministic Process Intelligence Engine (DTEAM) was engineered from the ground up to solve this exact problem. DTEAM abandons heuristic guessing in favor of a mathematically closed, axiomatically verifiable transformation kernel. Our governing principle is:
\begin{equation}
A = \mu(O^*)
\end{equation}
Where:
\begin{itemize}
    \item $O^*$ represents the closed ontology of the target process language (e.g., block-structured Workflow Nets).
    \item $\mu$ is a strictly deterministic transformation kernel (the RL engine + WASM runtime + reward shaping).
    \item $A$ is the resulting structural artifact (the Petri Net and its classification outcome).
\end{itemize}

By proving that $\mu$ is strictly deterministic, and by enforcing structural soundness mathematically within $O^*$, DTEAM ensures that $A$ is both optimal and verifiably free from overfitting.
